\begin{figure}[t]
\centering
\begin{tikzpicture}[>=Stealth,node distance=1.0cm,
 state/.style={draw,rounded corners=2pt,minimum width=1.45cm,minimum height=.85cm}]
\node[state] (a) {$L_1$};
\node[state,right=of a] (b) {$M_1$};
\node[state,right=of b] (c) {$M_1$};
\node[state,right=of c] (d) {$L_1$};
\node[state,right=of d] (e) {$L_{-1}$};
\draw[->] (a)--node[above,font=\small] {$A$} (b);
\draw[->] (b)--node[above,font=\small] {$I$} (c);
\draw[->] (c)--node[above,font=\small] {$A^{-1}$} (d);
\draw[->] (d)--node[above,font=\small] {$I$} (e);
\node[font=\small,align=center,below=0.5cm of c]
 {$I:t\mapsto-t$,\qquad $M_1(y,-t)=M_1(y,t)$\\[3pt]
 $L_0$ and the chosen finite collection of far endpoint jets return unchanged.};
\end{tikzpicture}
\caption{The continuation operation in the no-patch argument.
The alternate transverse branch is even at glancing. For each prescribed
finite set of far indices, the path \(A\) is chosen to fix both signs of
those indices. The operation produces another root of the same endpoint
polynomial; it requires no common path for infinitely many indices.}
\label{fig:monodromy}
\end{figure}
